\section{Domain templates}
\label{sec:domains}

The generic kernel is useful only if it leads to recognizable mathematics.  This section gives reusable
interfaces for the major domains represented in the repository.  The purpose is not to reprove every
packet's theorem, but to show exactly how its finite content enters the staged graph and where nonuniform
mathematical data remains.

\subsection{Exact finite algebra}

Finite algebra is the easiest domain because its operations are exact.  Integers use sign-and-magnitude or
paired-natural codes.  Rationals use normalized triples \((s,a,b)\) with \(b>0\) and
\(\gcd(a,b)=1\).  Polynomial codes contain a degree bound and a coefficient vector.  Matrices contain fixed
or list-coded dimensions and entry arrays.

A typical exact-algebra node has:
\begin{itemize}
  \item a validity formula for every code;
  \item a concrete \(\PR\) operation, such as addition, multiplication, determinant, substitution, or
        polynomial remainder;
  \item an equation theorem characterizing the operation;
  \item size bounds for its output code;
  \item a local correctness proof usable by specialization.
\end{itemize}

For polynomial identities, the strongest certificate is often a finite coefficient identity.  An assertion
\(f=g\) is reduced to equality of all coefficients below a known degree bound.  For matrix identities,
entrywise equality below the supplied dimensions is Booleanized.  Determinants can be defined by a fixed
permutation sum for fixed dimension or by a fraction-free elimination trace for variable dimension.  The
latter should carry pivot and row-operation certificates, so no unguarded division is used.

\paragraph{Zero-dimensional polynomial systems.}
Suppose a packet supplies a nonzero resultant or multiplication-matrix certificate showing that an affine
system has finitely many solutions, together with coordinate height bounds.  The framework enumerates all
rational coordinate tuples inside those bounds and verifies the equations.  The output list is complete by
the height theorem.  If only a dimension or cardinality bound is known, enumeration does not follow; the
missing height is an explicit boundary datum.

\paragraph{Finite quotient algebras.}
A quotient algebra may be represented by a basis, multiplication matrices, and proofs of the defining
relations.  Every polynomial expression reduces to finite matrix arithmetic.  Invertibility is guarded by
a determinant or a supplied inverse matrix.  Trace, norm, characteristic polynomial, and finite-dimensional
linear-algebra consequences are deterministic contracts.

\subsection{Represented real and complex analysis}

A quantitative real name is a primitive-recursive rational function \(a(q)\) with a displayed modulus, for
example
\[
 |a(q)-a(r)|\le 2^{-q}+2^{-r}.
\]
A complex name is a pair of real names.  The public type is not a second-order function object inside iPRA;
it is a fixed \(\PR\) symbol or a finite code for one member of a generated family, plus formulas describing
its values.

\paragraph{Arithmetic.}
Addition and negation have direct demand transformers.  Multiplication needs magnitude envelopes:
if \(|x|\le M_x\), \(|y|\le M_y\), and approximations have errors \(\delta_x,\delta_y\), then
\[
 |xy-\tilde x\tilde y|
 \le M_x\delta_y+M_y\delta_x+\delta_x\delta_y.
\]
Choose rational internal errors by a primitive-recursive rule so that the right side is at most the requested
output error.  Clear denominators to obtain integer inequalities suitable for Booleanization.

\paragraph{Reciprocal and division.}
A reciprocal is not total on names.  Supply an apartness guard \(|x|\ge g>0\).  If
\(|x-\tilde x|\le g/2\), then \(|\tilde x|\ge g/2\), and
\[
 \left|\frac1x-\frac1{\tilde x}\right|
 =\frac{|x-\tilde x|}{|x\tilde x|}
 \le \frac{2|x-\tilde x|}{g^2}.
\]
The demand transformer chooses input precision from \(g\) and the output request.  The exact operation is
therefore a gap contract.  Without \(g\), a relational interval avoiding no branch is the appropriate
output.

\paragraph{Limits and series.}
An infinite series theorem must provide a primitive-recursive tail modulus
\(N(q,x)\) satisfying a finite inequality
\[
 \left|\sum_{k=m}^{n}a_k(x)\right|\le 2^{-q}
 \quad(n\ge m\ge N(q,x)).
\]
The name at precision \(q\) is a finite partial sum at \(N(q+c,x)\).  Algebraic identities between series
are proved first for finite sums; tail bounds then transfer them to names.  Saying merely that a series
converges is insufficient for a producer contract, though a metastable finite-window theorem may still be
possible.

\paragraph{Differentiation and integration.}
A derivative contract consists of a local remainder bound, not an ideal limit:
\[
 |f(x+h)-f(x)-hD(x)|\le |h|\varepsilon(q).
\]
after a computed input demand.  An integral uses a finite quadrature rule and a modulus controlling the
mesh error.  The fundamental theorem of calculus becomes a composition of finite partition identities,
uniform continuity demands, and rational error budgets.  Quantifier-free induction handles partition sums.

\subsection{Formal and truncated power series}

Many generating-function arguments are finitary once stated coefficientwise.  A series at cutoff \(N\) is
a coefficient vector \((a_0,\ldots,a_{N-1})\).  Addition and Cauchy multiplication are exact finite folds.
Composition is guarded by a zero constant term or an explicit truncation discipline.  Inversion is guarded
by invertibility of the constant coefficient.

A theorem such as Lagrange inversion should first be stated as
\[
 [X^n]F(G(X))=R_n(a_0,\ldots,a_n)
\]
for arbitrary finite \(n\), where both sides are finite polynomial expressions.  Quantifier-free induction
on \(n\) or a finite combinatorial certificate proves the identity.  Analytic convergence radii, when
needed, are a separate represented-analysis layer with explicit tail bounds.  This separation prevents
unnecessary appeal to completed infinite series for a theorem whose actual consumer inspects finitely many
coefficients.

\subsection{Polynomial roots and spectral data}

For exact coefficients, a root certificate can be an isolating interval, disk, Sturm sequence, Rouch\'e
count, or multiplication-matrix certificate.  A complete root output is a finite cloud of regions with
multiplicities.  Refinement edges express containment and multiplicity preservation.

For matrices, approximate spectra can be represented by rational intervals or disks plus residual and
completeness certificates.  Ordered eigenvalues are stable for Hermitian matrices as a multiset and, with a
quantitative perturbation theorem, as sorted approximate numbers.  Eigenvectors and spectral projectors
need a gap when an individual eigenspace is requested.  At a collision, the invariant object is the combined
spectral subspace or projector for a separated cluster.

The relational-cloud framework therefore supplies a uniform vocabulary:
\begin{itemize}
  \item candidate cluster decompositions;
  \item compatibility under refinement;
  \item exact multiplicity and trace observables;
  \item diameter-stable root or spectral sets;
  \item exact labels only under separation.
\end{itemize}

\subsection{Finite combinatorics and probability}

Finite combinatorial objects---graphs, partitions, tableaux, permutations, finite groups, and finite
probability spaces---have direct list or bit-vector codes.  Counting identities are exact bounded sums.
A probabilistic statement in iPRA is ordinarily a rational inequality on a finite distribution or on a
finite approximation with explicit error.

For asymptotic probability, the theorem must state a modulus.  A central-limit approximation, for example,
is a family of finite inequalities
\[
 \sup_{t\in G_q}
 |F_n(t)-\Phi_q(t)|\le 2^{-q}
 \quad(n\ge N(q)),
\]
where \(G_q\) is a finite rational grid and \(\Phi_q\) is a finite approximation to the Gaussian
cumulative distribution.  Interpolation and tail bounds transfer the grid estimate to the requested finite
observation.  If a proof supplies only pointwise convergence without a modulus, the framework records it as
metastable or conditional rather than inventing \(N\).

\subsection{Finite topology}

General topological spaces are not first-order finite objects, but many formalizable theorems use finite
presentations:
\begin{itemize}
  \item finite simplicial or cubical complexes;
  \item finite posets and Alexandrov spaces;
  \item rational polyhedra with triangulations;
  \item finite open covers with an explicit nerve;
  \item mesh approximations with a modulus.
\end{itemize}
Boundary matrices, chain groups, incidence, subdivision, and simplicial maps are exact finite algebra.
Homology computations are proof-replayed matrix reductions.  A topological theorem becomes a theorem about
finite carriers plus approximation or invariance certificates.

For a compact metric object represented by a name, a finite net at precision \(q\) must be produced with a
coverage proof.  Compactness in ordinary prose does not automatically provide this net in iPRA.  If a net
producer and modulus are supplied, minimization of a uniformly continuous function is a finite search
followed by a quantitative transfer theorem.

\subsection{Finite categories and diagrams}

A finite category code consists of finite object and arrow sets, source and target maps, identities,
composition, and Booleanized category-law checks.  A finite functor is a pair of finite maps with law
certificates.  Natural transformations, limits of a fixed finite diagram, adjunction data, and finite Kan
extensions become bounded algebraic structures with exact verifiers.

The framework supports two kinds of uniformity:
\begin{enumerate}
  \item a theorem for all valid finite category codes, implemented by one list-coded primitive-recursive
        program;
  \item a generated theorem for each fixed finite presentation, in which the host compiler emits all
        tables directly.
\end{enumerate}
It does not quantify over a completed category of all mathematical structures.  Yoneda-style statements are
proved for coded hom-sets and finite diagrams, or for a fixed generated syntax whose operations have
concrete \(\PR\) symbols.

\subsection{Affine algebraic geometry and finite descent}

A finite affine chart is represented by polynomial equations, localization data, and bounded algebraic
certificates.  A principal open \(D(f)\) uses a witness to invert \(f\), usually an auxiliary variable
\(u\) with equation \(uf=1\).  Equality in a localization is witnessed by a finite exponent:
\[
 f^N(a-b)=0.
\]
A finite principal cover has a unit-ideal certificate
\[
 \sum_i h_i f_i^{N_i}=1.
\]
This certificate is exact and verifies coverage algebraically.

Gluing sections on a finite cover becomes a finite equalizer calculation.  Local sections have codes;
overlap equality has exponent witnesses; a finite Cech-style compatibility table is checked; and an
explicit algebraic gluing formula produces a global section.  The entire construction is a static proof
graph of finite algebra nodes and bounded searches.

Projective objects use homogeneous coordinate tuples with normalization or equivalence witnesses.  Heights
are explicit integer bounds on primitive coordinates.  Descent arguments are finite when supplied with a
finite cover of local cases, norm or resultant certificates, and height bounds.  The framework does not
supply those deep bounds automatically.

\subsection{Diophantine geometry}

The framework's role in Diophantine geometry is best summarized by the chain
\[
\begin{aligned}
 \text{global producer data}
 &\Longrightarrow \text{height cutoff or finite envelope}\\
 &\Longrightarrow \text{complete bounded enumeration}
 \Longrightarrow \text{verified list of points}.
\end{aligned}
\]
The last two implications are generic theorems.  The first is the substance of the deep mathematics.
For Runge-type curves, elementary coefficient estimates may give the height producer.  For a particular
elliptic curve, descent and finite covering data may give it.  For general Siegel or Faltings, the repository
reports correctly identify missing uniform effective inputs.  The staged kernel allows all downstream work
to be completed without claiming those inputs have been solved.

\subsection{Analytic number theory}

Finite sums over primes, divisor functions, and explicit sieve weights are directly primitive recursive.
A theorem such as a prime-counting estimate must be stated with explicit numerical constants and a threshold
\(N\).  If the analytic argument supplies them effectively, the finite inequality is an iPRA target.  If a
classical proof gives existence of a constant without a computable bound, the theorem is boundary-
conditional on a code for the constant and a finite verifier of the estimates it must satisfy.

This treatment is not a weakness of the framework.  It distinguishes the constructive content of an
explicit estimate from the nonuniform content of an ineffective existence proof.  It also permits partial
results: finite identities, explicit ranges, or metastable estimates can be fully embedded while the global
constant remains separate.

\subsection{A uniform application recipe}

Across all domains, a new application should proceed as follows.
\begin{enumerate}
  \item Choose exact finite codes and validity predicates.
  \item State the consumer's finite observation and demand.
  \item Separate producer, verifier, and soundness derivation.
  \item Replace unstable exact outputs by relations or add a gap.
  \item Compute backward demands and explicit resources.
  \item Prove local deterministic, loop, search, or relational contracts.
  \item Generate an arity-normalized static proof graph.
  \item Compile to \(\Prf\,\PRAAx\), then export to the native calculus if desired.
\end{enumerate}
The appendix gives a detailed audit checklist for this recipe.
